\chapter{Coalitions, Committees, and the Algebra of Aggregation}\label{ch:coalitions}

\epigraph{\itshape ``In every act of association, however small, the
problem of the relation of the part to the whole is already posed in
miniature.''}{Mancur Olson \cite{Olson1965}}

The previous chapter introduced the bare idea algebra
$\IdMon$. Its objects are individual ideas and its operation is the
binary composition $\compose$ that fuses any two ideas into a third.
A great deal of what passes for ``social'' or ``collective'' thought,
however, does not arise from binary acts at all but from the
\emph{simultaneous} contribution of indefinitely many parties: a
committee meeting, a research school, a jury, an actor-network, a
diffusion cascade. To talk about such phenomena inside the formalism
we need a free construction that turns finite multisets, or rather
finite \emph{sequences}, of ideas into a single idea while remaining
faithful to the monoidal structure underneath. That construction is
the \emph{coalition}, and the operator that carries it back into
$\Ideas$ is the \emph{aggregate}.

This chapter develops the algebra of coalitions and the two
\emph{structured} variants we shall need throughout the rest of the
book---\emph{committees} (a coalition with a distinguished chair) and
\emph{assemblies} (a coalition of committees). It proves the three
headline theorems of the chapter---coalition associativity, silent-agent
invariance, and the social congruence of $\sim$---and harvests the
corollaries that will be invoked, often without comment, in
Chapters~\ref{ch:grading} and onwards. We close, as in every chapter
of this monograph, with a six-entry zoo: Granovetter's strength of
weak ties, Olson's logic of collective action, Latour's
actor-networks, Bourdieu's distinction, Sawyer's emergent group
cognition, and Condorcet's jury theorem, each restated in the
vocabulary of the formalism and matched against the theorems just
proven.

The Lean witness for the development is the module
\texttt{IdeaTheory.Theorems10}; we cite individual results by their
informal names but reproduce no Lean code in the body, in keeping
with the conventions of the monograph.

\section{From single ideas to coalitions: the list monad over an idea algebra}
\label{sec:single-to-coalitions}

\subsection*{The free construction}

Let $\IdMon$ be a bare idea algebra in the sense of
Chapter~\ref{ch:foundations}: a set $\Ideas$ equipped with an
associative composition $\compose$ and a two-sided identity $\unit$.
Write $\List(\Ideas)$ for the set of finite sequences (or lists)
drawn from $\Ideas$. We use the standard list constructors: the empty
list $[\,]$, the cons $a :: \xs$ that prepends an element $a$ to a
list $\xs$, and the concatenation $\xs \mathbin{+\!\!+} \ys$. Two lists are
equal in $\List(\Ideas)$ exactly when they have the same length and
agree pointwise; in particular order matters.

\begin{definition}[Coalition]\label{def:coalition}
A \emph{coalition} of agents over the carrier $\Ideas$ is an element
\[
  \mathbf{c} = [a_1, a_2, \ldots, a_n] \in \List(\Ideas).
\]
The integer $n \geq 0$ is the \emph{size} of the coalition; the
elements $a_i \in \Ideas$ are the \emph{individually-held ideas} of
the agents. The empty coalition is denoted $[\,]$.
\end{definition}

The choice of \emph{lists} rather than \emph{multisets} or \emph{sets}
is deliberate. Following Arrow's discussion of the order in which
voters are polled \cite{Arrow1951} and Sen's parallel treatment in
\cite{Sen1970}, we treat the order of a coalition as a piece of
\emph{enumeration}---a contingent way of writing the coalition
down---rather than as a structural priority among its members. This
distinction is paid for in subsequent theorems: order-independence is
a result we will \emph{prove} (a special case of Theorem~\ref{thm:assoc}),
not an axiom we will \emph{assume}.

\subsection*{The list monad and naturality}

Categorically, the assignment $\Ideas \mapsto \List(\Ideas)$ is the
underlying functor of the \emph{free monoid monad} on the category of
sets. Its unit is the singleton inclusion
$\eta_\Ideas : \Ideas \to \List(\Ideas)$, $a \mapsto [a]$, and its
multiplication is concatenation,
$\mu_\Ideas : \List(\List(\Ideas)) \to \List(\Ideas)$, gluing a list
of lists into a single list. We will not need the abstract
machinery, but two consequences of it bear stating.

\begin{proposition}[Naturality of coalition formation]\label{prop:nat}
For every map of carriers $f : \Ideas \to \Ideas'$, the induced map
$\List(f) : \List(\Ideas) \to \List(\Ideas')$, $[a_1,\ldots,a_n]
\mapsto [f(a_1),\ldots,f(a_n)]$, is a homomorphism of free monoids.
In particular, the formation of coalitions commutes with relabelling
of agents.
\end{proposition}

\begin{proof}
Direct from the definition: $\List(f)([\,]) = [\,]$ and
$\List(f)(\xs \mathbin{+\!\!+} \ys) = \List(f)(\xs) \mathbin{+\!\!+} \List(f)(\ys)$.
\end{proof}

This naturality is the formal underwriting of the practice, common in
sociological and economic modelling, of treating the same coalition
under different ``descriptions'' of its members
\cite{ListPettit2011, Hardin1982}. So long as the description is
applied componentwise, the coalitional structure is undisturbed.

\subsection*{Notation and conventions}

We shall write $\Coalition(\Ideas)$ for $\List(\Ideas)$ when we want
to emphasise the social reading, and $\List(\Ideas)$ when emphasising
the carrier-set construction. The two are interchangeable. We write
$|\xs|$ for the size of a coalition, $\xs[i]$ for its $i$-th member
(zero-indexed), and $\xs \mathbin{+\!\!+} \ys$ for concatenation as above.
A coalition is \emph{flat} if it carries no extra structure, in
contrast to committees and assemblies introduced in
\S\ref{sec:comm-ass}. The empty coalition $[\,]$ should not be
confused with the singleton $[\unit]$ containing a single silent
agent: as lists they are different, even though, by
Corollary~\ref{cor:trivial-empty} below, they are coalition-equivalent.
This distinction matters in any framework that needs to track
membership beyond aggregation, e.g.~for accountability or audit
purposes \cite{ListPettit2011}.


\label{sec:aggregate}

The point of building $\List(\Ideas)$ at all is to compute the
\emph{single} idea that a coalition stands for. The operator that
does this is fixed by the monoid structure of $\IdMon$.

\begin{definition}[Aggregate]\label{def:aggregate}
The \emph{aggregate} of a coalition is the function
$\aggregate : \List(\Ideas) \to \Ideas$ defined by the recurrence
\[
  \aggregate([\,]) \;=\; \unit,
  \qquad
  \aggregate(a :: \xs) \;=\; a \compose \aggregate(\xs).
\]
\end{definition}

Equivalently, $\aggregate$ is the unique monoid homomorphism
$\List(\Ideas) \to \IdMon$ that sends each singleton $[a]$ to $a$;
existence and uniqueness are the universal property of the free
monoid \cite{MacLane1998}.

\subsection*{Small coalitions}

The recurrence yields immediately:

\begin{lemma}[Small-coalition aggregates]\label{lem:small}
For every $a, b, c \in \Ideas$,
\begin{align*}
  \aggregate([\,])     &= \unit,\\
  \aggregate([a])      &= a,\\
  \aggregate([a,b])    &= a \compose b,\\
  \aggregate([a,b,c])  &= a \compose (b \compose c).
\end{align*}
\end{lemma}

\begin{proof}
The first line is the base of the recurrence. For the second,
$\aggregate([a]) = a \compose \aggregate([\,]) = a \compose \unit = a$
by the right identity. The third unfolds twice and applies the second
to $\aggregate([b]) = b$. The fourth iterates the third.
\end{proof}

Note that $\aggregate$ is the \emph{right-leaning} fold: a triple is
read as $a \compose (b \compose c)$, not as $(a \compose b) \compose
c$. The two readings agree because $\compose$ is associative, but the
right-leaning convention is what allows the recurrence
$\aggregate(a :: \xs) = a \compose \aggregate(\xs)$ to be a single
clean equation, and it is this convention that we use throughout.

\subsection*{The chain operator}

In Chapter~\ref{ch:foundations} we introduced the \emph{chain
operator} $\chain{\cdot}$ that takes a finite indexed family
$\{a_i\}_{i=1}^{n}$ and returns its iterated composition. The
aggregate is exactly the chain operator restricted to coalitions
considered as lists:
\[
  \aggregate([a_1, \ldots, a_n]) \;=\; \chain{a_1, a_2, \ldots, a_n}
  \;=\; a_1 \compose a_2 \compose \cdots \compose a_n,
\]
with the empty chain understood to be $\unit$. The point of singling
out $\aggregate$ as a separate piece of vocabulary is not to introduce
a new operation but to license the social/cognitive readings developed
in this chapter and the next: chains are syntactic, aggregates are
\emph{joint judgements} of an aggregate of contributors
\cite{ListPettit2011}.

\subsection*{The append-decomposition law}

The single most useful technical fact about $\aggregate$ is that it
respects concatenation.

\begin{lemma}[Append decomposition]\label{lem:append}
For all coalitions $\xs, \ys \in \List(\Ideas)$,
\[
  \aggregate(\xs \mathbin{+\!\!+} \ys) \;=\; \aggregate(\xs) \compose \aggregate(\ys).
\]
\end{lemma}

\begin{proof}
By induction on $\xs$. If $\xs = [\,]$ then $\xs \mathbin{+\!\!+} \ys = \ys$
and the right-hand side is $\unit \compose \aggregate(\ys) =
\aggregate(\ys)$ by the left identity. For the inductive step, let
$\xs = a :: \xs'$. Then $\xs \mathbin{+\!\!+} \ys = a :: (\xs' \mathbin{+\!\!+} \ys)$, so
\begin{align*}
  \aggregate(a :: (\xs' \mathbin{+\!\!+} \ys))
    &= a \compose \aggregate(\xs' \mathbin{+\!\!+} \ys) \\
    &= a \compose (\aggregate(\xs') \compose \aggregate(\ys))
       \quad\text{(induction)}\\
    &= (a \compose \aggregate(\xs')) \compose \aggregate(\ys)
       \quad\text{(associativity)} \\
    &= \aggregate(a :: \xs') \compose \aggregate(\ys),
\end{align*}
which is the desired equality.
\end{proof}

\begin{corollary}[Trivial-side identities]\label{cor:trivside}
For every coalition $\xs$, $\aggregate(\xs \mathbin{+\!\!+} [\,]) =
\aggregate(\xs) = \aggregate([\,] \mathbin{+\!\!+} \xs)$.
\end{corollary}

\begin{proof}
Both equalities follow from Lemma~\ref{lem:append} together with
$\aggregate([\,]) = \unit$ and the two identity laws.
\end{proof}

\section{Committees vs.~assemblies: the structural distinction}
\label{sec:comm-ass}

Real social bodies are rarely flat. The minimal structural
elaboration is to mark one member out as the \emph{chair}---the
agenda-setter, the speaker-in-chief, the institutionally
distinguished node \cite{Sen1970, ListPettit2011}. The next is to
group flat coalitions into a meta-coalition---a council of councils,
an assembly of committees. Both elaborations admit a clean formal
treatment that we will use repeatedly.

\begin{definition}[Committee]\label{def:committee}
A \emph{committee} is a pair $\Comm = (\chair, \members)$ with
$\chair \in \Ideas$ and $\members \in \List(\Ideas)$. The
\emph{collective idea} of $\Comm$ is
\[
  \Comm.\collective \;=\; \chair \compose \aggregate(\members).
\]
\end{definition}

\begin{definition}[Assembly]\label{def:assembly}
An \emph{assembly} is a list of committees, $\Asm = [\Comm_1,
\ldots, \Comm_k]$. Its \emph{collective idea} is the aggregate of
the collective ideas of its committees:
\[
  \Asm.\collective \;=\; \aggregate([\Comm_1.\collective, \ldots,
  \Comm_k.\collective]).
\]
\end{definition}

The committee construction repackages the chair-cons of a list as a
single object. The next lemma makes the repackaging precise.

\begin{lemma}[Committee equals chair-prepended aggregate]\label{lem:chair-eq}
For every committee $\Comm = (\chair, \members)$,
\[
  \Comm.\collective \;=\; \aggregate(\chair :: \members).
\]
\end{lemma}

\begin{proof}
By Definition~\ref{def:committee} the left-hand side is $\chair
\compose \aggregate(\members)$, which by Definition~\ref{def:aggregate}
is exactly $\aggregate(\chair :: \members)$.
\end{proof}

\begin{lemma}[Empty / silent / singleton committees]\label{lem:comm-edge}
\begin{enumerate}
\item If $\members = [\,]$ then $\Comm.\collective = \chair$.
\item If $\chair = \unit$ then $\Comm.\collective = \aggregate(\members)$.
\item If $\Comm = (\chair, [\,])$ then $\Comm.\collective = \chair$.
\end{enumerate}
\end{lemma}

\begin{proof}
(1) and (3) are the same statement: $\chair \compose \aggregate([\,])
= \chair \compose \unit = \chair$. (2) is $\unit \compose
\aggregate(\members) = \aggregate(\members)$.
\end{proof}

For assemblies the corresponding small-cases are immediate from
Lemma~\ref{lem:small} applied to the list of committee aggregates:

\begin{lemma}[Small assemblies]\label{lem:asm-small}
$([\,]).\collective = \unit$; $([\Comm]).\collective =
\Comm.\collective$; $([\Comm, \Comm']).\collective = \Comm.\collective
\compose \Comm'.\collective$.
\end{lemma}

\subsection*{Why the distinction matters}

Flat coalitions, committees, and assemblies are distinguished
\emph{structurally} but not \emph{aggregationally}: each has a
collective idea in $\Ideas$, but the road from members to that idea
takes a different form. In Chapter~\ref{ch:emergence} we will see
that the emergence cocycle $\emerg$ \emph{does} distinguish them, and
this is precisely the formal content of Sawyer's claim that group
cognition is not a flattening of individual contributions
\cite{Sawyer2005, SawyerEmergence2005}. For the purposes of the
present chapter, however, the assembly-collective is fully determined
by its sequence of committee-collectives, and the committee-collective
is fully determined by its chair-prepended member coalition. So every
theorem about $\aggregate$ is automatically a theorem about
committees and assemblies, via Lemmas~\ref{lem:chair-eq}
and~\ref{lem:asm-small}.

\section{Silence: idempotence at zero and the silent coalition}
\label{sec:silence}

A coalition has many internal degrees of freedom that the aggregate
cannot detect. The simplest is silence.

\begin{definition}[Silent agent]\label{def:silent}
An agent $a \in \Ideas$ is \emph{silent} iff $a = \unit$. A
coalition $\xs$ is \emph{all-silent} iff every member is silent.
\end{definition}

\begin{definition}[Trivial coalition]\label{def:trivial}
For each $n \in \NN$, the \emph{trivial coalition of length $n$} is
\[
  \mathbf{0}_n \;=\; \underbrace{[\unit, \unit, \ldots, \unit]}_{n
  \text{ copies}} \;=\; \List.\!\replicate(n, \unit).
\]
\end{definition}

The basic fact about silence is that $\unit$ is the absorbing element
of $\aggregate$ from either side and at any depth.

\begin{lemma}[Silent insertion at the head]\label{lem:cons-id}
$\aggregate(\unit :: \xs) = \aggregate(\xs)$.
\end{lemma}

\begin{proof}
$\aggregate(\unit :: \xs) = \unit \compose \aggregate(\xs) =
\aggregate(\xs)$ by the left identity.
\end{proof}

\begin{lemma}[Trivial collapse]\label{lem:trivial}
$\aggregate(\mathbf{0}_n) = \unit$ for every $n \geq 0$.
\end{lemma}

\begin{proof}
By induction on $n$. The base $n = 0$ is $\aggregate([\,]) = \unit$.
For the inductive step, $\mathbf{0}_{n+1} = \unit ::
\mathbf{0}_n$, so by Lemma~\ref{lem:cons-id} $\aggregate(\mathbf{0}_{n+1})
= \aggregate(\mathbf{0}_n) = \unit$ by the inductive hypothesis.
\end{proof}

\begin{lemma}[All-silent collapse]\label{lem:all-silent}
If $\xs$ is all-silent then $\aggregate(\xs) = \unit$.
\end{lemma}

\begin{proof}
Induction on $\xs$. If $\xs = [\,]$ then $\aggregate(\xs) = \unit$.
If $\xs = a :: \xs'$ with $a$ silent and $\xs'$ all-silent, then
$\aggregate(\xs) = a \compose \aggregate(\xs') = \unit \compose \unit
= \unit$ by the inductive hypothesis and Lemma~\ref{lem:cons-id}.
\end{proof}

\begin{proposition}[Bulk silent absorption]\label{prop:bulk}
For every $\xs \in \List(\Ideas)$ and every $n \geq 0$,
\[
  \aggregate(\mathbf{0}_n \mathbin{+\!\!+} \xs)
  \;=\; \aggregate(\xs)
  \;=\; \aggregate(\xs \mathbin{+\!\!+} \mathbf{0}_n).
\]
In particular, $\aggregate(\unit :: \xs \mathbin{+\!\!+} [\unit]) = \aggregate(\xs)$.
\end{proposition}

\begin{proof}
By Lemma~\ref{lem:append}, $\aggregate(\mathbf{0}_n \mathbin{+\!\!+} \xs) =
\aggregate(\mathbf{0}_n) \compose \aggregate(\xs) = \unit \compose
\aggregate(\xs) = \aggregate(\xs)$, where the second equality is
Lemma~\ref{lem:trivial} and the third is the left identity. The
right-side identity is symmetric. The sandwich case follows by two
applications: $\aggregate(\unit :: \xs \mathbin{+\!\!+} [\unit]) =
\aggregate(\xs \mathbin{+\!\!+} [\unit]) = \aggregate(\xs)$ by
Lemma~\ref{lem:cons-id} and the right-side identity for $n=1$.
\end{proof}

The interpretive consequence will be drawn out in
\S\ref{sec:zoo} in connection with Olson's logic of collective action:
silent agents are formally invisible, and any model in which they are
not is implicitly using more structure than the bare aggregate provides.

\section{Coalition-equivalence: when do two coalitions aggregate to the same idea?}
\label{sec:equiv}

\begin{definition}[Coalition equivalence]\label{def:equiv}
Two coalitions $\xs, \ys \in \List(\Ideas)$ are
\emph{coalition-equivalent}, written $\xs \sim \ys$, iff
$\aggregate(\xs) = \aggregate(\ys)$.
\end{definition}

The relation $\sim$ is the kernel pair of the homomorphism
$\aggregate$. By a general fact about kernel pairs of monoid maps
\cite{Howie1995}, $\sim$ is an equivalence relation \emph{and} a
congruence with respect to concatenation.

\begin{proposition}[$\sim$ is an equivalence]\label{prop:equiv}
The relation $\sim$ on $\List(\Ideas)$ is reflexive, symmetric, and
transitive.
\end{proposition}

\begin{proof}
All three follow from the corresponding properties of equality in
$\Ideas$. Reflexivity: $\aggregate(\xs) = \aggregate(\xs)$.
Symmetry: equality is symmetric. Transitivity: equality is transitive.
\end{proof}

\begin{proposition}[$\sim$ is a congruence for $\mathbin{+\!\!+}$]\label{prop:cong}
If $\xs_1 \sim \xs_2$ and $\ys_1 \sim \ys_2$, then
$\xs_1 \mathbin{+\!\!+} \ys_1 \sim \xs_2 \mathbin{+\!\!+} \ys_2$. In particular,
$\sim$ is preserved by left-append and by right-append with any
fixed coalition.
\end{proposition}

\begin{proof}
By Lemma~\ref{lem:append} both sides expand:
\[
  \aggregate(\xs_1 \mathbin{+\!\!+} \ys_1) = \aggregate(\xs_1) \compose
  \aggregate(\ys_1) = \aggregate(\xs_2) \compose \aggregate(\ys_2) =
  \aggregate(\xs_2 \mathbin{+\!\!+} \ys_2),
\]
where the middle equality uses the two hypotheses.
\end{proof}

\begin{proposition}[Singleton equivalence is equality]\label{prop:single-equiv}
$[a] \sim [b] \iff a = b$.
\end{proposition}

\begin{proof}
By Lemma~\ref{lem:small}, $\aggregate([a]) = a$ and $\aggregate([b])
= b$, so $\aggregate([a]) = \aggregate([b]) \iff a = b$.
\end{proof}

\begin{corollary}[Trivial-empty equivalence]\label{cor:trivial-empty}
$\mathbf{0}_n \sim [\,]$ for every $n \geq 0$.
\end{corollary}

\begin{proof}
$\aggregate(\mathbf{0}_n) = \unit = \aggregate([\,])$ by
Lemma~\ref{lem:trivial}.
\end{proof}

\begin{remark}[$\sim$ vs.~list equality]
The quotient $\List(\Ideas)/\!\sim$ is, by the universal property of
the free monoid, isomorphic to the image of $\aggregate$ in
$\Ideas$. So coalition equivalence \emph{is} the difference between
\emph{which agents are present} (a list-level fact) and \emph{what
the coalition collectively says} (an idea-level fact). The whole
problem of social choice can be cast as the search for further
relations on $\List(\Ideas)$ that refine $\sim$ in interesting ways
\cite{Arrow1951, Sen1970}.
\end{remark}

\subsection*{Worked example: a five-member committee}

To fix ideas, take a committee with chair $c$ and four members $m_1,
m_2, m_3, m_4$, of whom the second is silent. Lemma~\ref{lem:chair-eq}
gives
\[
  \Comm.\collective \;=\; \aggregate([c, m_1, m_2, m_3, m_4]).
\]
Apply Lemma~\ref{lem:append} to split the list into a leading head
$[c, m_1]$ and a trailing tail $[m_2, m_3, m_4]$:
\[
  \Comm.\collective \;=\; (c \compose m_1) \compose \aggregate([m_2,
  m_3, m_4]).
\]
Now $m_2 = \unit$, so by Lemma~\ref{lem:cons-id} the trailing aggregate
collapses to $\aggregate([m_3, m_4]) = m_3 \compose m_4$, and the
final answer is $\Comm.\collective = c \compose m_1 \compose m_3
\compose m_4$. The silent member is invisible in the output even
though it is structurally present in the input. We will return to this
example when discussing the Olson zoo entry of \S\ref{sec:zoo}.

\subsection*{The quotient $\List(\Ideas)/\!\sim$}

It is worth recording the structural identity of the quotient
explicitly. Define $\overline{\aggregate} :
\List(\Ideas)/\!\sim\, \to \Ideas$ by $\overline{\aggregate}([\xs]) =
\aggregate(\xs)$. Proposition~\ref{prop:equiv} guarantees this is
well-defined; Lemma~\ref{lem:append} guarantees it is a monoid
homomorphism with respect to concatenation classes; and the universal
property of the free monoid identifies its image with the submonoid of
$\Ideas$ generated by the carriers of the coalitions in question. For
a generating set $S \subseteq \Ideas$, the image is the submonoid
$\langle S \rangle \subseteq \Ideas$. The quotient is therefore the
\emph{accessible part} of $\Ideas$ from the agents' inputs: a fact we
will use systematically when discussing the closed worlds of
Foucault-style epistemic systems in Chapter~\ref{ch:philosophy}.

\section{Insertion and delegation: the algebra of growing a coalition}
\label{sec:insertion}

A coalition is rarely born in its final form. It grows by accretion:
a new member joins at some position; a sub-coalition delegates
itself en bloc; a co-opted body is appended at the end. The
formalism mirrors these operations directly.

\begin{definition}[Positional insertion]\label{def:insert}
For $a \in \Ideas$, $n \in \NN$ and $\xs \in \List(\Ideas)$,
the \emph{insertion of $a$ at position $n$} is the list
$\Coalition.\insertAt(a, n, \xs)$ defined by
\[
  \insertAt(a, n, [\,]) = [a],
  \quad
  \insertAt(a, 0, \xs) = a :: \xs,
  \quad
  \insertAt(a, n{+}1, x :: \xs') = x :: \insertAt(a, n, \xs').
\]
If $n$ exceeds the length of $\xs$ the agent is appended at the end.
\end{definition}

\begin{definition}[Delegation]\label{def:delegate}
The \emph{delegation} of $\ys$ into $\xs$ is the concatenation
$\xs \mathbin{+\!\!+} \ys$.
\end{definition}

The decomposition law gives delegation an immediate algebraic interpretation.

\begin{lemma}[Delegation decomposition]\label{lem:delegate}
$\aggregate(\xs \mathbin{+\!\!+} \ys) = \aggregate(\xs) \compose
\aggregate(\ys)$. Delegation is associative as an operation on
coalitions: $(\xs \mathbin{+\!\!+} \ys) \mathbin{+\!\!+} \zs = \xs \mathbin{+\!\!+} (\ys
\mathbin{+\!\!+} \zs)$.
\end{lemma}

\begin{proof}
The decomposition is Lemma~\ref{lem:append}; the associativity is
the standard associativity of list concatenation \cite{Howie1995}.
\end{proof}

The algebraically important fact about insertion is that, when the
inserted agent is silent, insertion is the identity on aggregates
\emph{at every position}.

\begin{lemma}[Silent positional insertion]\label{lem:insert-id}
For every $n \geq 0$ and every $\xs \in \List(\Ideas)$,
\[
  \aggregate(\insertAt(\unit, n, \xs)) \;=\; \aggregate(\xs).
\]
\end{lemma}

\begin{proof}
Induction on $\xs$, with $n$ free. If $\xs = [\,]$, then
$\insertAt(\unit, n, [\,]) = [\unit]$ for all $n$ by
Definition~\ref{def:insert}, and $\aggregate([\unit]) = \unit \compose
\unit = \unit = \aggregate([\,])$.

If $\xs = x :: \xs'$ then we split on $n$. For $n = 0$,
$\insertAt(\unit, 0, x :: \xs') = \unit :: x :: \xs'$, and
$\aggregate(\unit :: x :: \xs') = \aggregate(x :: \xs')$ by
Lemma~\ref{lem:cons-id}. For $n = k + 1$, $\insertAt(\unit, k+1, x ::
\xs') = x :: \insertAt(\unit, k, \xs')$, and
\[
  \aggregate(x :: \insertAt(\unit, k, \xs'))
  = x \compose \aggregate(\insertAt(\unit, k, \xs'))
  = x \compose \aggregate(\xs')
  = \aggregate(x :: \xs'),
\]
where the middle equality is the inductive hypothesis applied at
position $k$.
\end{proof}

For non-silent insertions the situation is more subtle: in general
$\aggregate(\insertAt(a, n, \xs)) \neq \aggregate(\xs)$, and the
discrepancy is sensitive to where in the list $a$ has been placed.
This sensitivity is the formal source of the famous network-theoretic
asymmetries we will revisit in \S\ref{sec:zoo} under the headings of
Granovetter and Latour.

A useful midway lemma, which we will need for the headline theorems,
is that insertion at the head of a list whose head is itself silent
behaves like the swap of two silent agents:

\begin{lemma}[Head-swap-silent]\label{lem:swap-silent}
For every $a \in \Ideas$ and every $\xs \in \List(\Ideas)$,
$\aggregate(a :: \unit :: \xs) = \aggregate(\unit :: a :: \xs)$.
\end{lemma}

\begin{proof}
Both sides equal $a \compose \aggregate(\xs)$: the left by
$\aggregate(a :: \unit :: \xs) = a \compose (\unit \compose
\aggregate(\xs)) = a \compose \aggregate(\xs)$; the right by
$\aggregate(\unit :: a :: \xs) = \unit \compose (a \compose
\aggregate(\xs)) = a \compose \aggregate(\xs)$.
\end{proof}

\subsection*{Non-silent insertion: a structural lemma}

For a non-silent insertion the situation is governed by a simple
factorisation. Suppose $\xs = \xs_L \mathbin{+\!\!+} \xs_R$ is split at
position $n$. Then $\insertAt(a, n, \xs) = \xs_L \mathbin{+\!\!+} (a :: \xs_R)$,
and Lemma~\ref{lem:append} gives
\[
  \aggregate(\insertAt(a, n, \xs))
  \;=\; \aggregate(\xs_L) \compose a \compose \aggregate(\xs_R).
\]
Comparing with $\aggregate(\xs) = \aggregate(\xs_L) \compose
\aggregate(\xs_R)$, the discrepancy is the failure of $a$ to commute
out of the middle of the product. In particular, when $a$ is invertible
the two aggregates differ exactly by left-multiplication by the
conjugate $\aggregate(\xs_L) \cdot a \cdot \aggregate(\xs_L)^{-1}$;
this observation, made within Howie's general theory of semigroup
congruences \cite{Howie1995}, is the algebraic substance of the
Bourdieu zoo entry of \S\ref{sec:zoo}, where chair-position has a
preferred (head) location and a chair's contribution is read in a
distinguished frame.

\subsection*{Delegation chains}

A repeated delegation
$\xs_1 \mathbin{+\!\!+} \xs_2 \mathbin{+\!\!+} \cdots \mathbin{+\!\!+} \xs_k$ aggregates, by
$k-1$ applications of Lemma~\ref{lem:append}, to
\[
  \aggregate(\xs_1) \compose \aggregate(\xs_2) \compose \cdots
  \compose \aggregate(\xs_k).
\]
Two delegation chains with the same multiset of sub-aggregates are
$\sim$-equivalent in any commutative idea algebra; in the
non-commutative case they are $\sim$-equivalent only when they agree
as ordered tuples of $\compose$-factors. The contrast between these
two cases is precisely the contrast between Granovetter's
order-insensitive sociometric measures and the order-sensitive
``translation chains'' of Latour's actor-networks; we will return to
this in \S\ref{sec:zoo}.

\section{The headline theorems}
\label{sec:headline}

Three theorems organise the rest of the development. The first is
the social-composition law that says coalitions can be partitioned
arbitrarily before aggregation. The second is the
silent-agent invariance theorem that says abstainers are formally
invisible. The third is the social-congruence theorem that says
coalition equivalence is preserved by every coalitional move
introduced in this chapter.

\subsection*{Theorem 1: coalition associativity}

\begin{theorem}[Coalition associativity, Theorem~10.1]\label{thm:assoc}
For every $\xs, \ys, \zs \in \List(\Ideas)$, the aggregate of a
binary delegation decomposes as
\begin{equation}\label{eq:dec}
  \aggregate(\xs \mathbin{+\!\!+} \ys) = \aggregate(\xs) \compose \aggregate(\ys),
\end{equation}
and the aggregate of a triple delegation is independent of the
parenthesisation:
\begin{equation}\label{eq:re}
  \aggregate((\xs \mathbin{+\!\!+} \ys) \mathbin{+\!\!+} \zs)
  \;=\; \aggregate(\xs \mathbin{+\!\!+} (\ys \mathbin{+\!\!+} \zs)).
\end{equation}
\end{theorem}

\begin{proof}
Equation~\eqref{eq:dec} is exactly Lemma~\ref{lem:append}.
For \eqref{eq:re} we apply \eqref{eq:dec} on each side. The
left-hand side becomes
\[
  \aggregate((\xs \mathbin{+\!\!+} \ys) \mathbin{+\!\!+} \zs)
  = \aggregate(\xs \mathbin{+\!\!+} \ys) \compose \aggregate(\zs)
  = (\aggregate(\xs) \compose \aggregate(\ys)) \compose \aggregate(\zs).
\]
The right-hand side becomes
\[
  \aggregate(\xs \mathbin{+\!\!+} (\ys \mathbin{+\!\!+} \zs))
  = \aggregate(\xs) \compose \aggregate(\ys \mathbin{+\!\!+} \zs)
  = \aggregate(\xs) \compose (\aggregate(\ys) \compose \aggregate(\zs)).
\]
The two right-hand sides are equal by the associativity of $\compose$.
\end{proof}

The informal social-choice literature usually states the result for
two sub-coalitions only: ``the merger of two coalitions is the
$\compose$ of their two judgements'' \cite{Arrow1951, Sen1970}. The
formal version is sharper. It packages the two-coalition decomposition
\emph{and} the three-coalition reassociation into a single
proposition; it locates the only structural input as the
associativity of $\compose$; and it is, by~\eqref{eq:dec}, the
universal binary law from which every $n$-fold reassociation follows
by induction.

\subsection*{Theorem 2: silent-agent invariance}

\begin{theorem}[Silent-agent invariance, Theorem~10.2]\label{thm:silent}
For every coalition $\xs \in \List(\Ideas)$ the following four
identities hold:
\begin{enumerate}
\item $\aggregate(\insertAt(\unit, n, \xs)) = \aggregate(\xs)$ for
  every $n \geq 0$;
\item $\aggregate(\mathbf{0}_n \mathbin{+\!\!+} \xs) = \aggregate(\xs)$ for every
  $n \geq 0$;
\item $\aggregate(\xs \mathbin{+\!\!+} \mathbf{0}_n) = \aggregate(\xs)$ for every
  $n \geq 0$;
\item $\aggregate(\unit :: \xs \mathbin{+\!\!+} [\unit]) = \aggregate(\xs)$.
\end{enumerate}
\end{theorem}

\begin{proof}
Item (1) is Lemma~\ref{lem:insert-id}. For (2) and (3) apply the
append-decomposition law (Lemma~\ref{lem:append}) and then
Lemma~\ref{lem:trivial}: $\aggregate(\mathbf{0}_n \mathbin{+\!\!+} \xs) = \unit
\compose \aggregate(\xs) = \aggregate(\xs)$, with the symmetric
argument for (3). Item (4) is the sandwich case of
Proposition~\ref{prop:bulk}.
\end{proof}

The informal version, traceable to Sen's discussion of
``abstention-invariance'' in \cite{Sen1970}, treats only the
addition of one silent agent at the head. The formal statement is a
\emph{conjunction} of four invariances and is therefore strictly
sharper, in the precise sense that any one of (1)--(4) implies the
informal version but not conversely.

\subsection*{Theorem 3: social congruence}

\begin{theorem}[Social congruence, Theorem~10.3]\label{thm:congr}
The relation $\sim$ on $\List(\Ideas)$ enjoys the following seven
properties:
\begin{enumerate}
\item (Reflexivity) $\xs \sim \xs$ for every $\xs$.
\item (Symmetry) $\xs \sim \ys$ implies $\ys \sim \xs$.
\item (Transitivity) $\xs \sim \ys$ and $\ys \sim \zs$ imply $\xs \sim \zs$.
\item (Bilateral congruence) $\xs_1 \sim \xs_2$ and $\ys_1 \sim
  \ys_2$ imply $\xs_1 \mathbin{+\!\!+} \ys_1 \sim \xs_2 \mathbin{+\!\!+} \ys_2$.
\item (Left congruence) $\xs \sim \ys$ implies $\zs \mathbin{+\!\!+} \xs \sim
  \zs \mathbin{+\!\!+} \ys$ for every $\zs$.
\item (Right congruence) $\xs \sim \ys$ implies $\xs \mathbin{+\!\!+} \zs \sim
  \ys \mathbin{+\!\!+} \zs$ for every $\zs$.
\item (Silent insertion) $\xs \sim \insertAt(\unit, n, \xs)$ for every
  $n$ and every $\xs$.
\end{enumerate}
\end{theorem}

\begin{proof}
Items (1)--(3) are Proposition~\ref{prop:equiv}.
Item (4) is Proposition~\ref{prop:cong}. Items (5) and (6) are the
two specialisations of (4) obtained by taking $\zs \sim \zs$ as one
of the hypotheses, available by reflexivity. Item (7) is the
content of Lemma~\ref{lem:insert-id}: by definition of $\sim$,
$\xs \sim \insertAt(\unit, n, \xs)$ iff $\aggregate(\xs) =
\aggregate(\insertAt(\unit, n, \xs))$, which is exactly the lemma.
\end{proof}

The seven items together state that $\sim$ is, in the language of
\cite{Howie1995}, a monoid congruence on the free monoid
$(\List(\Ideas), \mathbin{+\!\!+}, [\,])$ that is moreover stable under the
silent-insertion move. This is the strongest statement that holds
\emph{without} commutativity. (For commutative idea algebras a
fourth structural axiom---closure under permutation of members---is
available; we discuss this in connection with Condorcet in
\S\ref{sec:zoo}.)

\subsection*{Corollaries}

The headline theorems immediately yield five corollaries that we
will need throughout the rest of the monograph. We collect them in a
single block.

\begin{corollary}[Two-coalition decomposition]\label{cor:two-decomp}
$\aggregate(\xs \mathbin{+\!\!+} \ys) = \aggregate(\xs) \compose \aggregate(\ys)$.
\end{corollary}

\begin{corollary}[Trivial-coalition collapse]\label{cor:trivial-collapse}
If every member of $\xs$ is silent, $\aggregate(\xs) = \unit$.
\end{corollary}

\begin{corollary}[Insertion-deletion congruence]\label{cor:insert-cong}
For every $n$ and every $\xs$, $\xs \sim \insertAt(\unit, n, \xs)$.
\end{corollary}

\begin{corollary}[Committee-collective decomposition]\label{cor:comm-decomp}
For every committee $\Comm = (\chair, \members)$,
$\Comm.\collective = \chair \compose \aggregate(\members)$.
\end{corollary}

\begin{corollary}[Assembly-pair decomposition]\label{cor:asm-decomp}
For any two committees $\Comm$, $\Comm'$,
$([\Comm, \Comm']).\collective = \Comm.\collective \compose
\Comm'.\collective$.
\end{corollary}

\begin{proof}
Each follows from the previous machinery:
Corollaries~\ref{cor:two-decomp}--\ref{cor:insert-cong} are restatements
of items in Theorems~\ref{thm:assoc} and~\ref{thm:silent};
Corollary~\ref{cor:comm-decomp} is the definition of $\Comm.\collective$;
Corollary~\ref{cor:asm-decomp} is Lemma~\ref{lem:asm-small} together
with Corollary~\ref{cor:two-decomp}.
\end{proof}

\subsection*{What the theorems do not say}

It is worth ending this section by stating what these results
\emph{do not} establish, and which structural augmentations would be
required to do so.

The aggregate is sensitive to order in general. Permuting the members
of a coalition changes the result unless $\compose$ happens to be
commutative on the members involved. So
``order-independence of joint judgement'' is \emph{not} a theorem of
the bare framework: it requires the additional axiom of commutativity,
which we will examine in connection with Condorcet's jury theorem in
\S\ref{sec:zoo}.

The aggregate of a non-trivial coalition with one inserted member is
\emph{not} in general $\compose$-related to the aggregate of the
original. There are coalition-equivalent pairs that differ by an
inserted non-silent member: e.g.~$[a, a^{-1}]$ in any group is
equivalent to $[\,]$. Cancellation, when present, supplies a far
richer theory of equivalence; this too is taken up below.

Finally, the aggregate is blind to the cocycle $\emerg$. We will see
in Chapter~\ref{ch:emergence} that the bare aggregate corresponds to a
\emph{trivialisation} of the cocycle---the value-of-aggregate is the
sum of values plus a correction governed by $\emerg$---and that the
correction is what makes Sawyer-style emergent group cognition
possible \cite{Sawyer2005}. The bare framework of this chapter
suffices, however, for everything in the zoo entries that follow.

\section{The zoo}\label{sec:zoo}

We close, as ever, with the chapter's six historical positions
matched against the formalism. Each entry restates a classical claim
in the vocabulary of $\aggregate$, $\compose$, $\sim$, and
silence, then uses the headline theorems to assess its formal
content.

\subsection*{Granovetter --- The strength of weak ties}

\begin{zooentry*}\thinker{Mark Granovetter (1973, 1983).}\end{zooentry*}

\textbf{Original claim.} In ``The strength of weak ties''
\cite{Granovetter1973}, Granovetter argues that the diffusion of
information through a social network is dominated by the rare,
non-redundant ``weak'' ties bridging dense cliques, rather than by
the many strong ties internal to those cliques. The thesis was
sharpened in \cite{GranovetterStrength1983} as the claim that the
expected informational reach of a network is essentially a function
of its bridges and not its hubs.

\formalclaim
Cast a person's network position as a coalition $\xs = [a_1, \ldots,
a_n]$ in which the agent's $i$-th tie is recorded as the contributed
idea $a_i \in \Ideas$. A \emph{strong} tie is one whose
contribution is silent or near-silent in the sense of
Definition~\ref{def:silent}: it returns nothing the agent does not
already have. A \emph{weak} tie is one whose contribution is
non-trivial: $a_i \neq \unit$. By
Theorem~\ref{thm:silent}, the silent (strong) ties may be inserted
or deleted at any position without changing $\aggregate(\xs)$.
Hence
\[
  \aggregate(\xs) \;=\; \aggregate(\xs_{\text{weak}}),
\]
where $\xs_{\text{weak}}$ is the sublist of non-silent contributions:
the coalition's aggregate is dominated, in the strict sense of
equality, by its non-$\unit$ insertions.

\formalstatus
True as a theorem of the bare framework. The required machinery is
Theorem~\ref{thm:silent}(1) applied repeatedly to the silent
contributions, together with Lemma~\ref{lem:append}.

\litcomment
The translation flattens an originally probabilistic claim about
\emph{expected} informational reach into a deterministic one about
\emph{actual} aggregates, but the structural insight is preserved
exactly: only the non-trivial insertions matter. The formalism also
supplies an immediate quantitative refinement, since the same
argument applied to the resonance pairing $\rho$ shows that the
\emph{strength} of an aggregate---measured by $\rho(\aggregate(\xs),
\aggregate(\xs))$---is unaffected by silent ties as well. This
matches the prediction of \cite{GranovetterStrength1983} that
sociometric measures restricted to ``effective'' edges should be
strictly less noisy than those that count all edges. A further
prediction unique to the formalism is that the bridging-versus-hub
asymmetry in informational reach is, at the level of the aggregate,
\emph{quantised}: silent contributions add nothing at all, rather than
merely a small amount, in any model that takes silence as an exact
identity. Empirical departures from this quantised limit are then
diagnostic of the failure of the silent-equals-$\unit$ idealisation
and motivate the resonance refinement of Chapter~\ref{ch:resonance}.

\subsection*{Olson --- The logic of collective action}

\begin{zooentry*}\thinker{Mancur Olson (1965).}\end{zooentry*}

\textbf{Original claim.} \emph{The Logic of Collective Action}
\cite{Olson1965}, sharpened in \cite{Hardin1982}, argues that in a
group facing the production of a public good, members below a
critical threshold of incentive will free-ride and contribute
nothing. The collective output is therefore the work of the
non-free-riders alone; the silent majority is, despite its size,
formally absent from the outcome.

\formalclaim
Model a producing group as a coalition $\xs = [a_1, \ldots, a_n]$
where each $a_i$ is the agent's contribution to the public-good
idea. A free-rider is an agent with $a_i = \unit$, i.e.~a silent
agent in the sense of Definition~\ref{def:silent}. By
Corollary~\ref{cor:trivial-collapse}, every all-silent
sub-coalition aggregates to $\unit$, and by Lemma~\ref{lem:append},
\[
  \aggregate(\xs) = \aggregate(\xs_{\text{contrib}}) \compose
  \unit = \aggregate(\xs_{\text{contrib}}),
\]
where $\xs_{\text{contrib}}$ is the sub-coalition of non-silent
agents. The free-riders contribute exactly $\unit$.

\formalstatus
True as a corollary of Theorems~\ref{thm:assoc}
and~\ref{thm:silent}.

\litcomment
The result is, formally, the dual of the Granovetter zoo entry: the
same theorem (silent-agent invariance) supplies both. What
distinguishes them is the interpretation of $\unit$: in Granovetter
it is the redundancy of a strong tie's information; in Olson it is
the free-rider's deliberate non-contribution. The formalism captures
both within a single absorbing element. The familiar Olsonian
prediction that public-good provision should fall \emph{nonlinearly}
with group size cannot be read off the bare framework---it requires
the cocycle $\emerg$ of Chapter~\ref{ch:emergence} together with the
grading $\deg$ of Chapter~\ref{ch:grading}---but the underlying
``silent agents are invisible'' fact is exactly
Corollary~\ref{cor:trivial-collapse}.

\subsection*{Latour --- Actor-network theory}

\begin{zooentry*}\thinker{Bruno Latour (2005).}\end{zooentry*}

\textbf{Original claim.} In \emph{Reassembling the Social}
\cite{Latour2005, LatourReassembling2005}, anticipated in Callon's
``sociology of translation'' \cite{Callon1986}, Latour argues that
social action is the joint effect of \emph{heterogeneous} networks
mixing human and non-human actors (instruments, documents, animals,
software). The key methodological move is the refusal to privilege
human contributions a~priori; the analyst follows whatever entity
acts.

\formalclaim
Take $\Ideas$ to contain both human-borne and non-human-borne ideas:
$\Ideas = \Ideas_H \cup \Ideas_N$. An actor-network is then a
coalition $\xs \in \List(\Ideas)$ with no constraint on which
sub-carrier each $a_i$ comes from. The collective effect of the
network is $\aggregate(\xs)$. By Proposition~\ref{prop:nat}, any
relabelling that respects the human/non-human distinction acts
componentwise on $\xs$ and so commutes with $\aggregate$; in
particular, no analytic privilege attaches to the human entries by
the structure of the formalism alone.

\formalstatus
True as a structural property of the bare framework. The relevant
piece is the naturality of $\aggregate$ under relabelling
(Proposition~\ref{prop:nat}), together with the absence of any
structural distinction between $\Ideas_H$ and $\Ideas_N$ at the
level of the monoid.

\litcomment
The formalisation supplies a precise version of Latour's
``methodological symmetry'' principle: there is no place in the
algebra of $\aggregate$ for a sortal restriction on the carrier.
Where the formalism diverges from ANT is in the use of \emph{ordered}
lists rather than networks: the chapter's coalitions impose an
enumeration that ANT explicitly rejects. The two views can be
reconciled by passing to the quotient $\List(\Ideas)/\!\sim$
(Definition~\ref{def:equiv}) and---under additional commutativity
hypotheses---to the quotient by permutation, where the result is
exactly the multiset structure ANT prefers. Latour's empirical
predictions, e.g.~that the introduction of a single non-human entry
can rearrange the entire aggregate, correspond formally to the
non-silent-insertion sensitivity discussed at the end of
\S\ref{sec:headline}.

\subsection*{Bourdieu --- Distinction}

\begin{zooentry*}\thinker{Pierre Bourdieu (1984, 1990).}\end{zooentry*}

\textbf{Original claim.} \emph{Distinction} \cite{Bourdieu1984} and,
in more methodological vein, \emph{The Logic of Practice}
\cite{BourdieuLogic1990}, argue that taste, judgement, and the
classifications of practical reason emerge from a graded, structured
field whose hierarchies are reproduced through the embodied
dispositions Bourdieu calls \emph{habitus}. A judgement is not the
output of an undifferentiated mass but of an internally graded body
in which some positions weigh more than others.

\formalclaim
Take a committee in the sense of Definition~\ref{def:committee} as
the formal counterpart of a Bourdieusian \emph{graded} body: the
chair carries the weight of institutional position, the members
carry the weight of dispersed dispositions. The aggregate norm of
the committee is then
\[
  \Comm.\collective \;=\; \chair \compose \aggregate(\members).
\]
Habitus is the global pattern by which $\chair$ is selected from the
members and the field induces an ordering on $\members$ before
aggregation. Distinction is the inequality
$\Comm.\collective \neq (\chair', \members).\collective$ when the
chair is changed.

\formalstatus
Contingent. The bare framework underwrites the committee
construction (Definition~\ref{def:committee}) and the structural
inequality $\chair \compose x \neq \chair' \compose x$ in general.
Whether it underwrites Bourdieu's stronger empirical claim that
distinction is reproduced \emph{stably} requires the additional
axiom that the chair-selection map is a fixed point of a higher-order
dynamics on $\Ideas$, which in turn needs the cocycle and grading
machinery of Chapters~\ref{ch:emergence} and~\ref{ch:grading}.

\litcomment
The formalisation is faithful to the structural part of Bourdieu's
account---the committee \emph{is} a graded coalition---and to the
predicted asymmetry between chair and members. It is unfaithful to
the embodied, dispositional aspect of habitus, which the bare
algebra cannot represent. A natural next step, taken up
in the chapter on grading, is to add a degree map $\deg : \Ideas \to
G$ such that $\deg(\Comm.\collective) = \deg(\chair) +
\deg(\aggregate(\members))$; this would give Bourdieu's ``volume of
capital'' a direct algebraic referent.

\subsection*{Sawyer --- Emergent group cognition}

\begin{zooentry*}\thinker{R. Keith Sawyer (2005).}\end{zooentry*}

\textbf{Original claim.} \emph{Social Emergence}
\cite{Sawyer2005, SawyerEmergence2005} argues that group cognition
is irreducible to the cognition of its members: the collective
output of a group exhibits properties not present in any individual
contribution and not predictable from the unstructured sum of
contributions. Sawyer's empirical exemplar is improvisational
theatre, where the joint scene is undeniably the work of the
ensemble but cannot be decomposed into separate authorial moves.

\formalclaim
For a coalition $\xs = [a_1, \ldots, a_n]$ define the
\emph{emergence defect}
\[
  \emerg(\xs) \;=\; v(\aggregate(\xs)) - \sum_{i=1}^{n} v(a_i),
\]
where $v : \Ideas \to A$ is any valuation into an abelian group $A$.
By Corollary~\ref{cor:two-decomp} together with the cocycle identity
of Theorem~4 (Chapter~\ref{ch:foundations}), $\emerg$ collects, for a
binary delegation, exactly the cocycle value
\[
  \emerg(\xs \mathbin{+\!\!+} \ys) - \emerg(\xs) - \emerg(\ys) = \emerg(\aggregate(\xs),
  \aggregate(\ys)).
\]
For a committee, the decomposition reads $\emerg(\Comm) =
\emerg(\chair, \aggregate(\members)) + \emerg([\chair]) +
\emerg(\members)$. The non-vanishing of $\emerg$ on a committee
aggregate is the formal content of group-level emergence.

\formalstatus
True as a corollary of the cocycle theorem (Theorem~4 of
Chapter~\ref{ch:foundations}) applied to the committee construction
of Definition~\ref{def:committee}. Whether $\emerg$ is non-trivial
on a particular group is empirical, but the formalism guarantees
that, when it is, the non-triviality is concentrated on the cocycle
$\emerg$ and not on the bare aggregate.

\litcomment
This is the chapter's only zoo entry that strictly requires
\emph{more} structure than the bare framework: a valuation $v$ and
its associated cocycle $\emerg$. The structural prediction is that,
under the cocycle decomposition, the irreducibility Sawyer
diagnoses is not located in the aggregate operator (which is fully
determined by the monoid) but in the failure of $v$ to be a
homomorphism. This rewards the methodological investment of the
cocycle apparatus and supplies a quantitative diagnostic---the
cocycle's cohomology class---for the strength of group emergence in
a given case.

\subsection*{Condorcet --- The jury theorem}

\begin{zooentry*}\thinker{Marquis de Condorcet (1785).}\end{zooentry*}

\textbf{Original claim.} The \emph{Essai}
\cite{Condorcet1785}, recovered analytically in
\cite{BlackCondorcet1958} and probabilistically in
\cite{YoungCondorcet1988}, proves that, under independence and
better-than-chance individual competence, the probability of correct
majority decision tends to one as the size of the jury grows. The
limiting object is an \emph{idealised aggregator} that perfectly
recovers the truth from a sufficiently large coalition.

\formalclaim
Take an idea algebra $\IdMon$ with $\compose$ commutative and with
left-cancellation: $a \compose c = b \compose c$ implies $a = b$.
For a coalition $\xs = [a_1, \ldots, a_n]$ in which each $a_i$ is
either a ``correct'' idea $t \in \Ideas$ or its inverse $t^{-1}$,
the aggregate
\[
  \aggregate(\xs) = a_1 \compose a_2 \compose \cdots \compose a_n
  = t^{k - (n-k)} = t^{2k - n}
\]
where $k$ is the number of $t$-votes. Coalition equivalence in this
setting is sensitive to the \emph{count} alone, not to the order:
$\xs \sim \ys$ iff their net $t$-power agrees. As $n \to \infty$
with each $a_i = t$ independently with probability $p > 1/2$, the
law of large numbers gives $k/n \to p$ almost surely, hence
$2k - n \to +\infty$ and $\aggregate(\xs) \to t^{\infty}$ in any
filtration that supports the limit.

\formalstatus
True under the additional axioms of commutativity and cancellation.
The bare framework gives only the structural skeleton (commutativity
of $\compose$ implies that $\sim$ is a permutation-invariant
congruence, by extending Theorem~\ref{thm:congr}); the probabilistic
limit requires standard additional hypotheses external to the algebra.

\litcomment
The Condorcet entry shows the opposite face of the formalism from
Granovetter's: where the latter exploits silent insertion to delete
redundant ties, the former exploits cancellation to make
\emph{antithetic} contributions cancel exactly. The two together
delimit the algebraic envelope of any aggregation theorem
expressible in the bare framework: at one extreme, all dissent
collapses to silence; at the other, dissent and assent cancel out
along the orbit of an inverse. Outside that envelope---non-commutative
votes, non-cancellative carriers, distributions on $\Ideas$ that do
not support a law of large numbers---more structure is required.
The grading and resonance machinery of subsequent chapters supply it.

\section*{Chapter summary}

We have introduced the coalition $\Coalition(\Ideas) =
\List(\Ideas)$ as the free monoidal hull of an idea algebra and the
aggregate $\aggregate$ as the unique homomorphism extending the
identity on singletons. We have proved the three headline
theorems---coalition associativity, silent-agent invariance, and
social congruence of $\sim$---and recorded their five corollaries.
The two structured variants, committees and assemblies, were
introduced as carriers of the chair-distinguished and
hierarchically-organised aggregations, respectively, and reduced to
the flat case via Lemmas~\ref{lem:chair-eq} and~\ref{lem:asm-small}.
The zoo entries cast Granovetter, Olson, Latour, Bourdieu, Sawyer,
and Condorcet within this single algebraic frame, locating each
position by which combination of bare-framework results, structural
augmentations, and external probabilistic hypotheses is required to
license its claims.

A useful way to read the chapter retrospectively is as a single
extended argument that the apparently informal social-choice notion of
``aggregate judgement'' is fully determined by three structural inputs
and nothing else: the monoid law of $\compose$, the identity element
$\unit$, and the free-list construction. Every result above is
extracted from these inputs by elementary inductive proofs. The
six historically motivated zoo entries then play a triangulation
role: each of them, on the customary informal reading, predicates of
``the group'' a property that the bare formalism either delivers
exactly (Granovetter, Olson), delivers in a methodologically-symmetric
form (Latour), delivers contingently on additional structural axioms
(Bourdieu), delivers only after enriching the framework with a
cocycle (Sawyer), or delivers under a substantive commutativity
hypothesis (Condorcet). The classification is a useful diagnostic
both for the formalism, in showing how much it can carry, and for the
historical positions, in distinguishing those structural commitments
that need only the bare monoid from those that smuggle in further
algebra.

The development of this chapter is the foundation for everything
that follows. The grading of Chapter~\ref{ch:grading} grades the
aggregate; the cocycle of Chapter~\ref{ch:emergence} measures the
defect of $v \circ \aggregate$ from a homomorphism; the resonance
pairing of Chapter~\ref{ch:resonance} extends to a bilinear form on
the coalition algebra. In each case the present chapter's three
theorems and five corollaries are the silent infrastructure on which
the new structure is laid. They will be invoked, often without
comment, throughout.
